\section{A matrix path with prescribed endpoints}
\label{sec:paths}

The construction in this section adapts the connecting-path argument of
\cite[Proposition~5.1]{ScalarCompanion} to a rectangular matrix $A$.
The action estimate uses a path from $K$ to $-K$ and penalizes failure
of $M_\delta(B)=AA^{\mathsf T}$, with $A=CU$ and $C$ lower triangular.
For a penalty parameter $\varepsilon>0$, we prescribe the difference
$M_\delta(B)-AA^{\mathsf T}$ as $\varepsilon B'$ through an ordinary
differential equation and choose $C$ to obtain the second endpoint.
The choice is made by continuation and boundary parity.
At the initial parameter there is exactly one solution after identifying
matrices that differ by row signs.  If no terminal solution existed,
smoothing and a small perturbation respecting these signs would produce
a quotient zero set that is a finite disjoint union of circles and
intervals, with exactly one boundary point.  Circles have no boundary
and intervals contribute two boundary points, giving a contradiction.

As in Section~\ref{sec:action-framework}, $\mathcal L$ is the vector
space of real lower triangular $n\times n$ matrices.  Set
$\kappa=\max_i k_i$.

\begin{proposition}[Connecting paths]
\label{prop:paths}
Let $n,d\geq1$, let $\Omega=(x_-,x_+)$ be a bounded interval, and let
$K=\operatorname{diag}(k_1,\ldots,k_n)$ with $k_i>0$.  Suppose
$U\in H^1_0(\Omega;\R^{n\times d})$ satisfies
\[
    \int_\Omega UU^{\mathsf T}\,dx=I_n.
\]
For every $\delta>0$ and $\varepsilon>0$, there are
$C_\varepsilon\in\mathcal L$ and
$B_\varepsilon\in C^1(\overline\Omega;\Sym_n)$ such that, with
$A_\varepsilon=C_\varepsilon U$,
\begin{equation}
\label{eq:path-ode}
    \varepsilon B_\varepsilon'
       =M_\delta(B_\varepsilon)-A_\varepsilon A_\varepsilon^{\mathsf T},
    \qquad
    B_\varepsilon(x_-)=K,
    \qquad B_\varepsilon(x_+)=-K.
\end{equation}
These choices satisfy
\begin{equation}
\label{eq:path-bounds}
    \sup_{x\in\overline\Omega}\|B_\varepsilon(x)\|_{\mathrm{op}}
       \leq\kappa,
    \qquad
    \sup_{0<\varepsilon\leq1}\|C_\varepsilon\|_{\HS}<\infty,
\end{equation}
where the second bound is for fixed $\delta$, $K$, and $\Omega$.
\end{proposition}

\begin{proof}
Fix $\delta>0$.  We use the continuous representative of $U$, available
because the interval is bounded and $U\in H^1$.  Choose a smooth
nonnegative function $\chi$ on $\Sym_n$, of compact support, which is
one on $\{B:\|B\|_{\mathrm{op}}\leq\kappa\}$ and is invariant under
conjugation by diagonal sign matrices.  Such a cutoff can be obtained
from any cutoff of a neighborhood of this compact set by averaging over
the finitely many sign matrices.  By Lemma~\ref{lem:field-regularity},
the field
\[
    \widehat M_\delta(B)=\chi(B)M_\delta(B)
\]
is bounded and globally Lipschitz.

\smallskip
\noindent\emph{The endpoint map.}
For fixed $\varepsilon>0$, $C\in\mathcal L$, and $0\leq\theta\leq1$,
solve the initial-value problem
\begin{equation}
\label{eq:path-homotopy}
    \varepsilon B'
       =\theta\widehat M_\delta(B)-(CU)(CU)^{\mathsf T},
    \qquad B(x_-)=K,
\end{equation}
and define
\[
    \Phi(C,\theta)=B(x_+)+K\in\Sym_n.
\]
The bounded Lipschitz field and the continuous forcing give a unique
$C^1$ solution on the whole interval.  Continuous dependence on the
parameters makes $\Phi$ continuous.  A zero of $\Phi(\cdot,1)$ gives
the required endpoints; we will then show that its path remains where
$\chi=1$.

Let $\mathcal G$ be the group of diagonal matrices whose entries are
$\pm1$.  It acts on $\mathcal L$ by $C\mapsto SC$ and on $\Sym_n$
by $B\mapsto SBS$.  The field is equivariant under these actions by
\eqref{eq:field-equivariance}, and $SKS=K$.  Uniqueness in
\eqref{eq:path-homotopy} therefore gives
\begin{equation}
\label{eq:path-equivariance}
    \Phi(SC,\theta)=S\Phi(C,\theta)S
    \qquad(S\in\mathcal G).
\end{equation}

All zeros of this endpoint map lie in a bounded set.  Indeed,
integrating \eqref{eq:path-homotopy} at a zero and using the
orthonormality of the rows of $U$ yields
\begin{equation}
\label{eq:path-integrated}
    CC^{\mathsf T}
       =2\varepsilon K
          +\theta\int_\Omega\widehat M_\delta(B(x))\,dx.
\end{equation}
Consequently, for $0<\varepsilon\leq1$,
\begin{equation}
\label{eq:path-C-bound}
    \|C\|_{\HS}^2
       \leq 2\tr K
          +n|\Omega|\sup_{B\in\Sym_n}
                    \|\widehat M_\delta(B)\|_{\mathrm{op}}.
\end{equation}
For each fixed $\varepsilon>0$, the same calculation gives a finite
bound.  Thus the joint zero set
$\{(C,\theta):\Phi(C,\theta)=0\}$ is compact.

The sign action is \emph{free} at every such zero: only the identity
sign matrix fixes $C$.  To prove this, it suffices to show that no row
of $C$ vanishes.  If row $i$ vanished, the reflection $S_i$ that changes
only its sign would fix $C$.  Equivariance and uniqueness would then
force $S_iB(x)S_i=B(x)$, so that $B(x)e_i=y(x)e_i$ for a scalar
function $y$.  The coordinate identity
\eqref{eq:field-coordinate} would give
\[
    \varepsilon y'
       =\theta\chi(B(x))\mu_\delta(k_i^2-y^2),
    \qquad y(x_-)=k_i.
\]
Here $\chi(B(x))$ is a fixed continuous coefficient.  Since
$\mu_\delta$ is locally Lipschitz and $\mu_\delta(0)=0$, uniqueness
for this scalar equation gives $y\equiv k_i$.  This contradicts the
terminal condition $y(x_+)=-k_i$.

At $\theta=0$ the endpoint map is explicit:
\begin{equation}
\label{eq:path-initial}
    \Phi(C,0)=2K-\varepsilon^{-1}CC^{\mathsf T}.
\end{equation}
Its zeros in $\mathcal L$ are precisely
\[
    C=\operatorname{diag}(d_1,\ldots,d_n),
    \qquad d_i=\pm\sqrt{2\varepsilon k_i}.
\]
For example, the first row of $CC^{\mathsf T}=2\varepsilon K$
determines $d_1$ and forces the remaining entries in the first column
of $C$ to vanish; induction gives the claim.  These $2^n$ zeros form
one orbit under $\mathcal G$.  Each is a regular zero, meaning that
the derivative with respect to $C$ is onto $\Sym_n$.  In fact, for a
lower triangular variation $H$, this derivative has diagonal entries
$-2\varepsilon^{-1}d_iH_{ii}$ and entries
$-\varepsilon^{-1}d_jH_{ij}$ when $i>j$.  It is therefore an
isomorphism between vector spaces of dimension $n(n+1)/2$.

\smallskip
\noindent\emph{Continuation after identifying row signs.}
We have a compact zero set, lying entirely where the sign action is
free, and a single regular zero at the initial parameter after taking
the quotient by that action.  We now show that these properties force
a zero at $\theta=1$.  The details below replace the continuous
endpoint map by a smooth one, so no differentiability of the ODE flow
with respect to $C$ is required.

Suppose, to the contrary, that $\Phi(C,1)$ never vanishes.  Compactness
allows us to choose a bounded invariant open set $O\subset\mathcal L$
containing the $C$-projection of every zero, with $\overline O$ contained
in the open set where all rows are nonzero.  There is then a positive
lower bound for $\|\Phi\|_{\HS}$ on
\begin{equation}
\label{eq:path-gap-set}
    (\partial O\times[0,1])\ \cup\ (\overline O\times\{1\}).
\end{equation}
Reparametrize $\theta$ by a smooth map $\rho:[0,1]\to[0,1]$ that
equals zero near $0$ and has $\rho(1)=1$.  The resulting map
$F(C,\theta)=\Phi(C,\rho(\theta))$ agrees with the smooth map
\eqref{eq:path-initial} on an initial collar, that is, for all
$0\leq\theta\leq a$ with some $a>0$.  It retains a positive gap on
\eqref{eq:path-gap-set}.

Approximate $F$ uniformly on $\overline O\times[0,1]$ by a smooth
equivariant map $\Psi$ that still equals \eqref{eq:path-initial} on a
smaller initial collar.  Here is a direct construction.  Extend $F$
constantly in $\theta$ outside $[0,1]$, smooth by convolution in the
finite-dimensional variables, and use a smooth function of $\theta$
to splice the approximation to $F(C,0)$ inside its initial collar.
Averaging the result in the form
\[
    \frac1{|\mathcal G|}\sum_{S\in\mathcal G}
       S\Psi_0(SC,\theta)S
\]
makes it equivariant.  Uniform approximation can be as close as
needed, so the positive gap persists.

We next make zero a regular value in the interior by a small
equivariant perturbation.  The freeness of the action on $\overline O$
is exactly what permits this step.  Around any $C_0\in\overline O$,
choose a ball whose translates by distinct sign matrices are disjoint,
and a smooth bump function $b$ supported in that ball with $b(C_0)=1$.
For $H\in\Sym_n$, the map
\[
    V_H(C)=\sum_{S\in\mathcal G} b(SC)SHS
\]
is smooth and equivariant, and $V_H(C_0)=H$.  Taking $H$ in a basis of
$\Sym_n$ gives maps whose values span the target near $C_0$.
A finite covering of $\overline O$ therefore supplies smooth
equivariant maps $V_1,\ldots,V_N$ whose values span $\Sym_n$ at every
point of $\overline O$.  The action on the target need not be free:
the disjoint supports in the domain allow an arbitrary target value
at a representative of each orbit.

Choose a smooth function $h:[0,1]\to[0,\infty)$ that vanishes near
$0$ and is positive thereafter, with its zero set contained in the
collar where $\Psi(C,\theta)=\Phi(C,0)$.  For a parameter
$\alpha=(\alpha_1,\ldots,\alpha_N)\in\R^N$, put
\[
    \Psi_\alpha(C,\theta)
       =\Psi(C,\theta)+h(\theta)\sum_{j=1}^N\alpha_jV_j(C).
\]
Restrict $\alpha$ to a sufficiently small open ball about zero so that
the gap on \eqref{eq:path-gap-set} remains.  At an interior zero of the
map of all three variables $(C,\theta,\alpha)$, its differential is
onto: if $h(\theta)>0$, the parameter directions span the target;
if $h(\theta)=0$, the derivative in $C$ is the isomorphism already
computed for \eqref{eq:path-initial}.

The interior zero set of this parameter family is thus a smooth
manifold.  Apply Sard's theorem \cite{Sard1942} to its projection onto the parameter
ball.  For a regular value $\alpha$ of this projection, zero is a
regular value of $(C,\theta)\mapsto\Psi_\alpha(C,\theta)$.  To see
the implication, write the differential of the parameter family as
$(L_1,L_2)$ in the $(C,\theta)$ and $\alpha$ directions.  Surjectivity
of the differential of the projection means that, for every $w$,
there is a $v$ with $L_1v+L_2w=0$.  Thus
$\operatorname{ran}L_2\subset\operatorname{ran}L_1$; since
$(L_1,L_2)$ is onto, $L_1$ is onto as well.  Sard's theorem supplies
such parameters arbitrarily close to zero.

For this choice of $\alpha$, the zero set of $\Psi_\alpha$ in
$O\times[0,1]$ is a compact smooth one-dimensional manifold with
boundary.  Its dimension follows from
$\dim\mathcal L=\dim\Sym_n$.  There are no zeros on
\eqref{eq:path-gap-set}, and near $\theta=0$ the map is the initial
map, independent of $\theta$, with invertible derivative in $C$.
The boundary therefore consists precisely of the $2^n$ initial zeros.

The sign group acts freely on this manifold.  Its quotient is again
a compact one-dimensional manifold with boundary: around each point,
choose a neighborhood disjoint from its nontrivial translates, so
that the quotient has the same local interval or half-interval
coordinate.  All initial zeros belong to one orbit, so the quotient
has exactly one boundary point.  This is impossible.  Every compact
one-dimensional manifold is a finite disjoint union of circles and
closed intervals, and hence has an even number of boundary points.
The contradiction proves that the original map $\Phi(\cdot,1)$ has
a zero.  This is the usual regular-value proof of invariance of parity,
applied after quotienting by row signs; see also
\cite[Sections~2--4]{Milnor1965}.

\smallskip
\noindent\emph{Confinement of the path.}
Let $C$ and $B$ be given by such a zero.  For any fixed real unit
vector $e$, set $y(x)=e^{\mathsf T}B(x)e$.  By
\eqref{eq:field-jensen}, whenever $|y(x)|>\kappa$,
\[
    e^{\mathsf T}M_\delta(B(x))e
       \leq\mu_\delta\!\left(e^{\mathsf T}K^2e-y(x)^2\right)<0.
\]
Since $\chi\geq0$ and $(CU)(CU)^{\mathsf T}\geq0$, the differential
equation gives $y'(x)\leq0$ on this set.  Both endpoint values of $y$
lie in $[-\kappa,\kappa]$.  A connected interval on which $y>\kappa$
cannot begin at the level $\kappa$ while $y$ is nonincreasing.
Likewise, an interval on which $y<-\kappa$ cannot end at the level
$-\kappa$ while $y$ is nonincreasing.  The first endpoint excludes
the former excursion, and the second endpoint excludes the latter.
Thus $|e^{\mathsf T}B(x)e|\leq\kappa$ for every unit vector $e$ and
every $x$.

It follows that $\|B(x)\|_{\mathrm{op}}\leq\kappa$ throughout the
interval, so $\chi(B(x))=1$ and \eqref{eq:path-homotopy} becomes
\eqref{eq:path-ode}.  Finally, \eqref{eq:path-C-bound} gives the
uniform bound on $C_\varepsilon$ in \eqref{eq:path-bounds} for fixed
$\delta$.  This completes the construction.
\end{proof}
